\chapter{Cadrage et cahier des charges}

\section{Objectifs métier, techniques et pédagogiques}

Pour bien démarrer le projet, j'ai structuré mes objectifs en trois catégories. Cette séparation m'aide à vérifier que je réponds aux besoins des joueurs, que je fais les bons choix techniques et que je valide les compétences attendues pour le titre CDA.

\subsection*{Objectifs métier}
\begin{center}
    \begin{tabular}{|p{4cm}|p{6cm}|p{3cm}|}
        \hline
        \textbf{Objectif} & \textbf{Pourquoi c'est important} & \textbf{Livrable} \\
        \hline
        Améliorer l'expérience des joueurs & Passer le taux d'abandon de 78\% à 30\% & Plateforme opérationnelle \\
        \hline
        Faire gagner du temps aux joueurs & 25 minutes gagnées par session = 833h/mois pour 500 utilisateurs & Module de planning intégré \\
        \hline
        Fidéliser la communauté & Augmenter de 25\% le temps sur les activités complexes & Système de badges \\
        \hline
        Centraliser les informations & Éviter aux joueurs de consulter 3-4 sources différentes & Guides unifiés \\
        \hline
    \end{tabular}
\end{center}

\subsection*{Objectifs techniques}
\begin{center}
    \begin{tabular}{|p{4cm}|p{6cm}|p{3cm}|}
        \hline
        \textbf{Objectif} & \textbf{Pourquoi c'est important} & \textbf{Livrable} \\
        \hline
        Temps de réponse < 2s & Améliorer l'expérience utilisateur et réduire les abandons & Monitoring New Relic \\
        \hline
        Gérer 100 joueurs simultanés & Tenir les pics après les mises à jour du jeu & Architecture scalable \\
        \hline
        Garantir 99\% de disponibilité & Un service de raid doit être fiable & Infrastructure redondante \\
        \hline
        Sécuriser les comptes & Protéger les données utilisateur (conformité RGPD) & Audit de sécurité \\
        \hline
        Maintenir le code & Couverture de tests > 80\% pour éviter la dette technique & Pipeline CI/CD \\
        \hline
    \end{tabular}
\end{center}

\subsection*{Objectifs pédagogiques}
\begin{center}
    \begin{tabular}{|p{4cm}|p{6cm}|p{3cm}|}
        \hline
        \textbf{Objectif} & \textbf{Lien avec le titre CDA} & \textbf{Livrable} \\
        \hline
        Maîtriser React/Node.js & Compétence cœur du développement applicatif & Code source documenté \\
        \hline
        Concevoir une architecture 3-tiers & Compétence en architecture logicielle & Diagrammes techniques \\
        \hline
        Intégrer une API complexe (Bungie) & Savoir connecter des services externes & Connecteur API fonctionnel \\
        \hline
        Mettre en place une stratégie de tests & Compétence en assurance qualité & Rapports de couverture \\
        \hline
        Déployer en continu (CI/CD) & Compétence en déploiement et maintenance & Pipeline opérationnel \\
        \hline
    \end{tabular}
\end{center}

\section{Comment je mesure la réussite du projet}

Pour savoir si le projet avance dans la bonne direction, j'ai défini des indicateurs précis. Je les suis régulièrement pendant le développement.

\begin{center}
\begin{tabular}{|p{3cm}|p{2.5cm}|p{2cm}|p{2cm}|p{2cm}|}
\hline
\textbf{Indicateur} & \textbf{Source} & \textbf{Fréquence} & \textbf{Résultat actuel} & \textbf{Objectif} \\
\hline
Temps d'affichage des pages & New Relic & Temps réel & 1.8s & < 2s \\
\hline
Temps pour planifier un raid & Logs utilisateur & Par session & 18 min & < 10 min \\
\hline
Connexions réussies & Logs backend & Quotidien & 92\% & > 95\% \\
\hline
Utilisateurs actifs mensuels & Google Analytics & Mensuel & 0 (lancement) & 500 \\
\hline
Couverture des tests & GitHub Actions & À chaque PR & 75\% & > 80\% \\
\hline
Abandon au premier raid & Tracking & Hebdomadaire & 78\% & < 30\% \\
\hline
Satisfaction utilisateur & Sondages & Mensuel & N/A & ≥ 4.5/5 \\
\hline
\end{tabular}
\end{center}

\subsection*{Exemple de suivi avec GitHub Projects}

Voici à quoi ressemble mon tableau de bord GitHub Projects. Il me permet de visualiser l'avancement en un coup d'œil :

\begin{verbatim}
Colonnes : Backlog → À faire → En cours → Revue → Terminé

Backlog (6 tâches) :
- Authentification OAuth Bungie [Priorité haute]
- Guides interactifs [Priorité haute]  
- Gestion des escouades [Priorité haute]
- Calendrier collaboratif [Priorité moyenne]
- Profils joueurs [Priorité moyenne]
- Système de badges [Priorité basse]

En cours (2 tâches) :
- Maquettes UI (85% terminé)
- Configuration environnement (90% terminé)

Terminé (3 tâches) :
- Spécifications fonctionnelles
- Architecture technique
- Étude de marché
\end{verbatim}

\subsection*{Validation avec de vrais joueurs}

Le 15 février 2026, j'ai organisé une session de validation avec trois profils types : Thomas (débutant), Sarah (leader expérimentée) et Alex (joueur très technique).

\textbf{Ce qu'ils m'ont appris :}
\begin{itemize}
    \item Thomas avait du mal avec certains termes techniques → j'ajoute un glossaire dans les guides.
    \item Sarah a trouvé la création d'escouade très intuitive, mais voulait inviter ses coéquipiers plus vite → je simplifie le processus à 3 clics maximum.
    \item Alex aimerait voir sa progression plus clairement → j'ajoute des indicateurs visuels dans les guides.
\end{itemize}

Tous ont validé les maquettes et le concept général. Je peux continuer en étant sûr de répondre à leurs vrais besoins.

\section{Pourquoi j'ai choisi cette stack technique}

\subsection*{PostgreSQL + Prisma + React/Node.js}

\textbf{PostgreSQL :}
\begin{itemize}
    \item Je peux définir des règles précises entre les tables (clés étrangères, contraintes d'unicité), ce qui garantit que les données des joueurs et des escouades restent cohérentes.
    \item Le moteur gère très bien les requêtes complexes, comme quand je dois chercher des joueurs selon plusieurs critères.
    \item Le type JSONB me donne de la flexibilité pour stocker des réglages ou des données de jeu qui peuvent varier.
    \item Les transactions ACID assurent que si je crée une escouade et que j'ajoute son leader, tout se fait en une seule fois, sans risque d'incohérence.
\end{itemize}

\textbf{Prisma :}
\begin{itemize}
    \item Il génère automatiquement les types TypeScript à partir de mon schéma de base. Je gagne du temps et j'évite des erreurs.
    \item Les migrations sont versionnées : je peux revenir en arrière si jamais une modification pose problème.
    \item L'autocomplétion dans mon éditeur me fait gagner un temps précieux au quotidien.
\end{itemize}

\subsection*{Pourquoi j'ai écarté d'autres solutions}

\begin{center}
\begin{tabular}{|p{3cm}|p{4cm}|p{4cm}|}
\hline
\textbf{Alternative} & \textbf{Avantages} & \textbf{Raison du rejet} \\
\hline
MySQL & Mature et performant & PostgreSQL gère mieux les données JSON, utiles pour moi \\
\hline
TypeORM & Populaire, multi-bases & Prisma offre une meilleure sécurité des types \\
\hline
SQLite & Simple à configurer & Ne passerait pas à l'échelle avec plusieurs utilisateurs \\
\hline
Firebase & Développement rapide & Je voulus éviter la dépendance à un fournisseur unique \\
\hline
\end{tabular}
\end{center}

\section{Priorisation des fonctionnalités (MoSCoW)}

J'ai classé les fonctionnalités en quatre catégories pour être sûr de livrer l'essentiel en premier.

\begin{center}
    \begin{tabular}{|p{2.5cm}|p{3.5cm}|p{4cm}|}
        \hline
        \textbf{Priorité} & \textbf{Fonctionnalité} & \textbf{Pourquoi c'est indispensable} \\
        \hline
        \textbf{Must Have} & Authentification Bungie & Sans ça, pas de compte utilisateur \\
        \textbf{(Essentiel)} & Guides interactifs & C'est la raison d'être du projet \\
        & Gestion des escouades & Le cœur de la collaboration \\
        & Base de données PostgreSQL & Rien ne fonctionne sans données \\
        \hline
        \textbf{Should Have} & Calendrier collaboratif & Pour le gain de temps mesurable \\
        \textbf{(Important)} & Profil + statistiques & Pour l'engagement des joueurs \\
        \hline
        \textbf{Could Have} & Système de badges & Pour la motivation, mais pas vital \\
        \textbf{(Accessoire)} & Notifications & Évite les oublis, mais on peut s'en passer au début \\
        \hline
        \textbf{Won't Have} & Application mobile native & Trop coûteux pour la V1 \\
        \textbf{(Pas pour maintenant)} & Chat en temps réel & Discord fait déjà ça très bien \\
        & Streaming vidéo & Hors sujet pour un MVP \\
        \hline
    \end{tabular}
\end{center}

\subsection*{Ce que contiendra la version 1.0 (MVP)}
\begin{itemize}
    \item Authentification complète avec Bungie
    \item Guides pour 3 raids : Vault of Glass, Last Wish, Deep Stone Crypt
    \item Création et gestion d'escouades
    \item Calendrier simple pour planifier des sessions
    \item Profil basique avec statistiques principales
\end{itemize}

\section{Scénarios de test concrets (format Gherkin)}

Pour chaque fonctionnalité, j'ai écrit des scénarios qui décrivent exactement comment elle doit se comporter. Ça m'aide à coder juste ce qu'il faut et à vérifier que ça marche.

\subsection*{Scénario 1 : Connexion avec Bungie}
\begin{verbatim}
Étant donné un joueur non connecté
Quand il clique sur "Se connecter avec Bungie"
Et qu'il s'authentifie sur le site de Bungie
Et qu'il autorise l'application
Alors il est redirigé vers son tableau de bord
Et son profil est créé ou mis à jour automatiquement
\end{verbatim}

\subsection*{Scénario 2 : Consultation d'un guide}
\begin{verbatim}
Étant donné un joueur connecté
Quand il choisit le raid "Vault of Glass"
Alors le guide s'affiche avec toutes les étapes
Et les mécaniques importantes sont illustrées
Et le temps estimé (45-60 min) est indiqué
Et les équipements recommandés sont listés
\end{verbatim}

\subsection*{Scénario 3 : Création d'une escouade}
\begin{verbatim}
Étant donné un joueur connecté
Quand il crée une nouvelle escouade "Les Raiders du Dimanche"
Et qu'il fixe la taille maximum à 6 joueurs
Et qu'il invite 5 amis par leur pseudo Bungie
Alors l'escouade apparaît dans sa liste
Et les invitations sont envoyées
Et il peut voir qui a déjà accepté
\end{verbatim}

\section{Les personnes concernées et les risques}

\subsection*{Les trois profils types d'utilisateurs}

\textbf{Thomas - Le débutant motivé (25 ans)}
\begin{itemize}
    \item \textbf{Profil :} Joue depuis 2 mois, 50 heures de jeu
    \item \textbf{Ce qui le frustre :} "Je comprends rien aux mécaniques des raids", "Personne veut de moi car je suis débutant", "Je perds une heure à chercher des infos"
    \item \textbf{Ce dont il a besoin :} Des guides clairs, une équipe patiente, un environnement où il peut apprendre sans stress
\end{itemize}

\textbf{Sarah - La leader expérimentée (30 ans)}
\begin{itemize}
    \item \textbf{Profil :} 2000+ heures de jeu, organise 3 raids par semaine
    \item \textbf{Ce qui la frustre :} "Je perds 30 minutes à tout organiser", "Je répète toujours les mêmes explications", "Les outils sont éparpillés"
    \item \textbf{Ce dont elle a besoin :} Centraliser ses outils, gagner du temps, suivre la progression de son équipe
\end{itemize}

\textbf{Alex - Le stratège (35 ans)}
\begin{itemize}
    \item \textbf{Profil :} Créateur de contenu, il adore optimiser les stratégies
    \item \textbf{Ce qui le frustre :} "Les infos ne sont jamais à jour", "Pas de données pour comparer les stratégies"
    \item \textbf{Ce dont il a besoin :} Des analytics fiables, des données en temps réel
\end{itemize}

\subsection*{Principaux risques et comment je les anticipe}

\begin{center}
    \begin{tabular}{|p{3.5cm}|p{4cm}|p{4cm}|}
        \hline
        \textbf{Risque} & \textbf{Impact} & \textbf{Comment je m'y prépare} \\
        \hline
        L'API Bungie change & Élevé & Je teste régulièrement et je surveille les annonces \\
        \hline
    Peu de joueurs adoptent l'outil & Élevé & J'implique des testeurs dès le début pour coller à leurs besoins \\
        \hline
        L'application rame & Moyen & Je fais des tests de charge tôt et j'optimise au fur et à mesure \\
        \hline
        Une faille de sécurité & Critique & Je fais auditer le code et je suis les bonnes pratiques OWASP \\
        \hline
    \end{tabular}
\end{center}

\section{Les fonctionnalités en détail}

\subsection*{Pour les joueurs (Front Office)}

\begin{center}
    \begin{tabular}{|p{3.5cm}|p{6cm}|}
        \hline
        \textbf{Fonctionnalité} & \textbf{Description} \\
        \hline
        Connexion Bungie & S'identifier avec son compte Bungie, gestion sécurisée de la session \\
        \hline
        Guides interactifs & Suivre les raids étape par étape, avec illustrations et conseils \\
        \hline
        Gestion d'escouade & Créer un groupe, inviter des joueurs, gérer les rôles \\
        \hline
        Calendrier & Planifier des sessions, voir les disponibilités de son équipe \\
        \hline
        Profil personnel & Voir ses statistiques, son historique, ses badges \\
        \hline
        Recherche de joueurs & Trouver des coéquipiers selon leur niveau, disponibilité, style de jeu \\
        \hline
    \end{tabular}
\end{center}

\subsection*{Pour l'administration (Back Office)}

\begin{center}
    \begin{tabular}{|p{3.5cm}|p{6cm}|}
        \hline
        \textbf{Fonctionnalité} & \textbf{Description} \\
        \hline
    Gestion des utilisateurs & Modérer les contenus, gérer les signalements \\
        \hline
        Gestion des guides & Créer, modifier, valider les guides de raids \\
        \hline
        Statistiques & Voir comment la plateforme est utilisée \\
        \hline
        Sauvegardes & Protéger les données avec des backups automatiques \\
        \hline
    \end{tabular}
\end{center}

\subsection*{Respect de la vie privée (RGPD)}

J'ai intégré les règles RGPD dès la conception :

\begin{center}
    \begin{tabular}{|p{3cm}|p{7cm}|}
        \hline
        \textbf{Point clé} & \textbf{Ce que j'ai mis en place} \\
        \hline
        Données collectées & Pseudo Bungie, stats de jeu, préférences \\
        \hline
    Consentement & L'utilisateur accepte explicitement en se connectant \\
        \hline
        Sécurité & Chiffrement des données sensibles, HTTPS partout \\
        \hline
        Durée de conservation & 3 ans après la dernière connexion, puis suppression \\
        \hline
        Droits des utilisateurs & Possibilité d'exporter ou supprimer ses données depuis son profil \\
        \hline
    \end{tabular}
\end{center}

\section{Feuille de route (Roadmap)}

\subsection*{Version 1.0 - Mars 2026 (MVP)}
\begin{itemize}
    \item On peut se connecter avec Bungie
    \item On consulte des guides pour 3 raids
    \item On crée une escouade et on invite des amis
    \item On planifie une session dans un calendrier simple
\end{itemize}

\subsection*{Version 1.1 - Juillet 2026}
\begin{itemize}
    \item Chat intégré pour discuter en direct avec son escouade
    \item Recommandations d'équipement personnalisées
    \item Interface améliorée selon les retours des premiers utilisateurs
    \item Version anglaise en plus du français
\end{itemize}

\subsection*{Version 1.2 - Septembre 2026}
\begin{itemize}
    \item Statistiques avancées pour les leaders
    \item Connexion avec Discord (webhooks)
    \item Guides pour tous les raids du jeu
    \item Objectif : 300 utilisateurs actifs par mois
\end{itemize}

\subsection*{Version 2.0 - Janvier 2027}
\begin{itemize}
    \item Application mobile (React Native)
    \item API ouverte pour que d'autres développeurs créent des outils
    \item Intégration de vidéos et de streams
    \item Objectif : 500+ utilisateurs actifs par mois
\end{itemize}